%==================================================================
% Ini adalah abstrak dalam bahasa indonesia 
%==================================================================

%% DILARANG EDIT BAGIAN INI
\clearpage
\phantomsection
\addcontentsline{toc}{chapter}{ABSTRAK}
\begin{center}
    \textbf{\large{\judulid}}\\[2em]
    \penulis\\
    NIM \nim\\[2em]
    \textbf{ABSTRAK}\\[0.5cm]
\end{center}
%% DILARANG EDIT BAGIAN INI

%% edit bagian ini
INFINITE adalah divisi di bawah UKM Rekayasa Teknologi (Restek) UNY yang berfungsi sebagai wadah pengembangan diri mahasiswa di bidang teknologi informasi. Sistem informasi organisasi INFINITE sebelumnya terfragmentasi dalam dua sistem, yaitu Laravel dengan arsitektur monolitik di \textit{shared hosting} dan CMS Strapi yang terpisah, sehingga menimbulkan tantangan pemeliharaan, konsistensi data, inefisiensi biaya operasional, serta ketidakmampuan integrasi dengan aplikasi eksternal dan infrastruktur Kubernetes. Penelitian ini bertujuan mengembangkan sistem informasi organisasi INFINITE yang terkontainerisasi dengan arsitektur \textit{headless} menggunakan Laravel sebagai \textit{backend}, Next.js sebagai \textit{frontend}, dan autentikasi berstandar protokol \textit{OpenID Connect} (OIDC), serta di-\textit{deploy} di Kubernetes dengan pendekatan GitOps menggunakan ArgoCD.

Pengembangan sistem dilakukan menggunakan metode \textit{Research and Development} (R\&D) dengan model \textit{Waterfall} melalui empat tahapan: \textit{Requirements Definition}, \textit{System and Software Design}, \textit{Implementation and Unit Testing}, serta \textit{Integration and System Testing}. Pengujian dilakukan berdasarkan standar ISO/IEC 25010:2023 untuk \textit{Product Quality} dan ISO/IEC 25019:2023 untuk \textit{Quality in Use}, yang mencakup aspek \textit{Functional Suitability}, \textit{Interaction Capability}, \textit{Performance Efficiency}, dan \textit{Beneficialness}.

Hasil penelitian menunjukkan bahwa sistem berhasil dikembangkan dan di-\textit{deploy} sebagai aplikasi terkontainerisasi di Kubernetes, dengan hasil pengujian \textit{Functional Suitability} sebesar 100\% (Sangat Layak), \textit{Interaction Capability} 91{,}73\% (Sangat Layak), \textit{Performance Efficiency} pada \textit{frontend} berupa rata-rata skor Lighthouse 90{,}42 (Baik) dan pada \textit{backend} berupa \textit{Response Time} 1 detik, \textit{Throughput} 53{,}6 permintaan per detik, \textit{Error Rate} 0\%, serta maksimum 189 \textit{Virtual Users} simultan, dan \textit{Beneficialness} 81{,}30\% (Sangat Layak). Berdasarkan hasil tersebut, sistem informasi organisasi INFINITE tergolong layak digunakan sebagai solusi pengelolaan data organisasi secara terintegrasi dan berkelanjutan.\\[0.6cm]
%% edit sampai sini

%% DILARANG EDIT BAGIAN INI
\noindent \textbf{Kata kunci:} Sistem Informasi, \textit{Headless}, Kontainer.
%% DILARANG EDIT BAGIAN INI